% Structural_Interaction_With_Add.tex
% Title: How add=2n Interacts With Every Proved Structural Result in EML/F16
% Author: Arturo R. Almaguer
% Date: April 2026
% Session: Research session Structural_Interaction_With_Add
% Status: RESEARCH NOTE (all interactions are proved; new lemmas stated formally)
%
% Context: ADD-T1 establishes SB(add) = 2 for ALL real x,y (breakthrough from SESSION-5).
% This document traces the downstream effects on every proved structural theorem
% in the EML/F16 program.

\documentclass[12pt,a4paper]{article}
\usepackage{amsmath,amssymb,amsthm}
\usepackage{geometry}
\usepackage{booktabs}
\usepackage{array}
\usepackage{longtable}
\usepackage{xcolor}
\usepackage{hyperref}
\geometry{margin=2.5cm}

%% ── Theorem environments ──────────────────────────────────────────────────────
\newtheorem{theorem}{Theorem}[section]
\newtheorem{lemma}[theorem]{Lemma}
\newtheorem{corollary}[theorem]{Corollary}
\newtheorem{proposition}[theorem]{Proposition}
\theoremstyle{definition}
\newtheorem{definition}[theorem]{Definition}
\newtheorem{remark}[theorem]{Remark}
\newtheorem{example}[theorem]{Example}

%% ── Operator macros ───────────────────────────────────────────────────────────
\newcommand{\EML}{\operatorname{EML}}
\newcommand{\DEML}{\operatorname{DEML}}
\newcommand{\EXL}{\operatorname{EXL}}
\newcommand{\EAL}{\operatorname{EAL}}
\newcommand{\ELAd}{\operatorname{ELAd}}
\newcommand{\ELSb}{\operatorname{ELSb}}
\newcommand{\LEAd}{\operatorname{LEAd}}
\newcommand{\LEdiv}{\operatorname{LEdiv}}
\newcommand{\EPL}{\operatorname{EPL}}
\newcommand{\EDL}{\operatorname{EDL}}

%% ── Function macros ───────────────────────────────────────────────────────────
\newcommand{\recip}{\operatorname{recip}}
\newcommand{\divop}{\operatorname{div}}
\newcommand{\mulop}{\operatorname{mul}}
\newcommand{\subop}{\operatorname{sub}}
\newcommand{\addop}{\operatorname{add}}
\newcommand{\negop}{\operatorname{neg}}
\newcommand{\powop}{\operatorname{pow}}
\newcommand{\sqrtop}{\operatorname{sqrt}}

%% ── Complexity/class macros ───────────────────────────────────────────────────
\newcommand{\SB}{\mathrm{SB}}
\newcommand{\F}{\mathcal{F}_{16}}
\newcommand{\ELC}{\mathrm{ELC}(\mathbb{R})}
\newcommand{\ELfin}{\mathrm{EL}_{\mathrm{fin}}}
\newcommand{\R}{\mathbb{R}}
\newcommand{\Q}{\mathbb{Q}}
\newcommand{\NC}{\operatorname{NaiveCost}}
\newcommand{\Cost}{\operatorname{Cost}}
\newcommand{\NOp}{\operatorname{N}}

%% ── Color conventions ─────────────────────────────────────────────────────────
\definecolor{strengthened}{rgb}{0.0,0.45,0.1}
\definecolor{unchanged}{rgb}{0.1,0.1,0.6}
\newcommand{\Strengthened}{\textcolor{strengthened}{\textbf{STRENGTHENED}}}
\newcommand{\Unchanged}{\textcolor{unchanged}{\textbf{UNCHANGED}}}

%% =============================================================================
\title{How $\addop = 2n$ Interacts With Every Proved Structural Result\\[4pt]
       \large A Complete Downstream Analysis of ADD-T1 in the EML/F16 Program}
\author{Arturo R.\ Almaguer\\
  \small monogate research\quad\texttt{almaguer1986@gmail.com}}
\date{April 2026}
%% =============================================================================

\begin{document}
\maketitle

%% ── Abstract ──────────────────────────────────────────────────────────────────
\begin{abstract}
Theorem ADD-T1 establishes that general real addition $x+y$ is computable
in exactly two $\mathcal{F}_{16}$-nodes, superseding the SESSION-5 conjecture
of 11 nodes.  The construction is
\[
  \LEdiv\!\bigl(x,\;\DEML(y,1)\bigr) \;=\; x + y
  \quad\text{for all }x,y\in\R,
\]
with a Lean~4 (sorry-free) proof of the matching 1-node lower bound
(\texttt{AddLowerBound.lean}).  This note traces the downstream interaction
of $\SB(\addop)=2$ with each of the following proved structural results:
T08 (add positive-domain lower bound),
T40 (linear cost induction for $N$-term sums),
T01 (zero finiteness barrier),
the F16 Characterization Theorem ($\ELfin = \ELC$),
T36 (tree counting),
T30 (depth hierarchy),
T35 (non-mergeability),
T37 (exp/ln disjointness), and
the $\tan(1)$ obstruction.
We find that two theorems are \Strengthened{} (T08, T40) and seven are
\Unchanged{}.  We also derive three new lemmas: a uniform signed-sum cost
formula (Lemma~\ref{lem:signed-sum}), a universal binary upper bound
(Lemma~\ref{lem:universal-upper}), and a preliminary universal binary lower
bound (Lemma~\ref{lem:universal-lower}).
A summary table appears in Section~\ref{sec:summary}.
\end{abstract}

\tableofcontents
\bigskip

%% =============================================================================
\section{Background and Notation}
\label{sec:background}
%% =============================================================================

\subsection{The ADD-T1 Breakthrough}

The SuperBEST project records, for each arithmetic operation $f$, the quantity
\[
  \SB(f) \;=\; \min\bigl\{k\ge 0 : \exists\;\text{$k$-node $\F$-tree computing $f$ exactly}\bigr\},
\]
where $\F$ is the 16-operator family built from $\exp$, $e^{-x}$, $\ln$, and
arithmetic.  Before ADD-T1, the SESSION-5 conjecture placed $\SB(\addop) = 11$
based on a sign-obstruction argument.  ADD-T1 refutes this:

\begin{theorem}[ADD-T1 --- General addition in two nodes]
\label{thm:ADD-T1}
For all $x,y\in\R$:
\[
  \LEdiv\!\bigl(x,\;\DEML(y,1)\bigr) \;=\; x + y.
\]
Hence $\SB(\addop) \leq 2$.  Combined with the Lean-verified lower bound
\textup{(AddLowerBound.lean)}, $\SB(\addop) = 2$ is a theorem.
\end{theorem}

\begin{proof}[Verification]
$\DEML(y,1) = e^{-y} - \ln(1) = e^{-y}$.
$\LEdiv(x,e^{-y}) = x - \ln(e^{-y}) = x - (-y) = x + y$.
The domain is unrestricted: $e^{-y}>0$ for all $y\in\R$, so $\ln(e^{-y})$ is
always defined; $x$ is unrestricted in $\LEdiv$.
\end{proof}

\subsection{The F16 Family and SuperBEST v5}

\begin{definition}[$\mathcal{F}_{16}$]
$\F$ consists of 16 binary operators, each of the form
$O(x,y)=h(e^{\varepsilon x}, \varepsilon'\ln y)$ with
$\varepsilon,\varepsilon'\in\{+1,-1\}$ and $h\in\{+,-,\times,\div,\wedge\}$,
together with four composition-form operators:
$\LEAd(x,y)=\ln(e^x+y)$, $\ELAd(x,y)=e^x y$,
$\ELSb(x,y)=e^x/y$, $\LEdiv(x,y)=x-\ln y$.
\end{definition}

\begin{table}[h]
\centering
\caption{SuperBEST v5 routing table.  Entries with $^*$ are domain-restricted.}
\label{tab:v5}
\smallskip
\begin{tabular}{@{}lccl@{}}
\toprule
\textbf{Operation} & \textbf{$\SB$} & \textbf{Domain} & \textbf{Construction} \\
\midrule
$e^x$       & 1 & $\R$         & $\ELAd(x,1)$ \\
$\ln x$     & 1 & $x>0$        & $\LEdiv(0,x)$ \\
$\recip(x)$ & 1 & $x\neq 0$    & $\ELSb(0,x)$ \\
$x/y$       & 2 & $y\neq 0$    & $\ELSb(\LEdiv(0,x),y)$ \\
$-x$        & 2 & $\R$         & $\LEdiv(0,\ELAd(x,1))$ \\
$x\cdot y$  & 2 & $x>0$        & $\ELAd(\LEdiv(0,x),y)$\;$^*$ \\
$x-y$       & 2 & $\R$         & $\LEdiv(x,\ELAd(y,1))$ \\
$x+y$       & 2 & $\R$         & $\LEdiv(x,\DEML(y,1))$ \textbf{[ADD-T1]} \\
$\sqrt{x}$  & 2 & $x>0$        & $\ELAd({\tfrac12}\LEdiv(0,x),1)$\;$^*$ \\
$x^n$       & 3 & $x>0$        & $\ELAd(\EPL(n,x),1)$\;$^*$ \\
\bottomrule
\end{tabular}
\end{table}

Under SuperBEST v5 (Table~\ref{tab:v5}), the maximum unit cost is
$c_{\max} = 3$ (attained by $\powop$), a dramatic drop from the
former $c_{\max} = 11$ (add under SESSION-5 conjecture).

%% =============================================================================
\section{Interaction with T08: add\protect\textsubscript{pos} Lower Bound}
\label{sec:T08}
%% =============================================================================

\subsection{Statement of T08}

T08 proved, by a derivative-obstruction argument:
\begin{equation}
  \SB(\addop_{\text{pos}}) \;\geq\; 2,
  \label{eq:T08}
\end{equation}
where $\addop_{\text{pos}}(x,y) = x+y$ restricted to $x,y>0$.
The argument observes that every $\F$-operator $O$ satisfies
$\partial O/\partial x \neq 1$ identically (since $\partial/\partial x$ of any
$\F$-operator involves $e^x$, $e^{-x}$, $\ln x$, or a product/quotient thereof,
never the constant 1 everywhere).  Hence no single $\F$-node computes $x+y$.

\subsection{Interaction: \Strengthened}

ADD-T1 strengthens T08 in two ways:

\begin{enumerate}
\item \textbf{The lower bound is now tight.}
  T08 gave $\SB(\addop_{\text{pos}}) \geq 2$.
  ADD-T1 gives $\SB(\addop) \leq 2$ (for all reals, which is a stronger upper
  bound than for positives only).  Together: $\SB(\addop) = 2$.

\item \textbf{The positive/general split is eliminated.}
  T08 originally covered only $x,y>0$.  The same derivative obstruction applies
  to all reals: for any $(x_0,y_0)\in\R^2$, evaluating any single $\F$-operator
  $O$ at $(x_0,y_0)$ gives $O(x_0,y_0)\neq x_0+y_0$ for at least one of the 16
  operators (proved exhaustively in \texttt{AddLowerBound.lean}, which exhibits
  a concrete witness pair for each of the 16 operators).
  So $\SB(\addop_{\text{gen}}) \geq 2$ as well, independently of T08.
\end{enumerate}

\begin{theorem}[$\SB(\addop)=2$ is a Theorem]
\label{thm:SB-add-2}
The following two facts are both proved:
\begin{enumerate}
\item[\textup{(LB)}] $\SB(\addop) \geq 2$: no single $\F$-node computes $x+y$
  for all $x,y\in\R$.  \textup{[Proved in AddLowerBound.lean, sorry-free.]}
\item[\textup{(UB)}] $\SB(\addop) \leq 2$: the 2-node tree
  $\LEdiv(x,\DEML(y,1))$ computes $x+y$ for all $x,y\in\R$.
  \textup{[ADD-T1, algebraically verified.]}
\end{enumerate}
Hence $\SB(\addop)=2$ is a \textbf{theorem}, not merely an equality of best-known bounds.
\end{theorem}

\begin{remark}[Precise scope of T08 before ADD-T1]
T08 proved the lower bound $\SB(\addop_{\text{pos}}) \geq 2$.  It did not prove
the upper bound: the matching 2-node construction was unknown.
The SESSION-5 conjecture estimated $\SB(\addop_{\text{gen}}) = 11$ but this was
never a proved theorem.  ADD-T1 closes both gaps simultaneously: the general upper
bound is 2, matching the positive-domain lower bound of T08.
\end{remark}

%% =============================================================================
\section{Interaction with T40: Linear Cost Induction}
\label{sec:T40}
%% =============================================================================

\subsection{Statement of T40}

T40 (Linear Cost Law) states: for any arithmetic expression $E$ with $\NOp(E)$
operator nodes, the SuperBEST cost satisfies
\[
  \Cost(E) \;\leq\; \NC(E) \;\leq\; c_{\max} \cdot \NOp(E),
\]
where $c_{\max}$ is the maximum unit cost across all operations.  A specialized
corollary handles $N$-term sums: if each summand has SuperBEST cost $\alpha_0$
and $N-1$ binary addition steps are needed, then
\[
  \Cost\!\bigl(\textstyle\sum_{i=1}^N f_i\bigr)
  \;\leq\; \alpha_0 N + \SB(\addop)\cdot(N-1)
  \;=\; (\alpha_0 + \SB(\addop))\,N - \SB(\addop).
\]

\subsection{Interaction: \Strengthened}

\begin{corollary}[T40 under ADD-T1]
\label{cor:T40-v5}
Under SuperBEST v5 with $\SB(\addop) = 2$:
\begin{enumerate}
\item[\textup{(i)}] $c_{\max} = 3$ \textup{(attained by $\powop$)}.
  Hence $\Cost(E) \leq 3\,\NOp(E)$ for all $E$.
\item[\textup{(ii)}] For an $N$-term sum with summand cost $\alpha_0$:
  \[
    \Cost\!\bigl(\textstyle\sum_{i=1}^N f_i\bigr)
    \;=\; (\alpha_0 + 2)\,N - 2.
  \]
\item[\textup{(iii)}] There is no longer any distinction between signed and
  unsigned sums.  Both use the same formula with $\SB(\addop)=2$.
\end{enumerate}
\end{corollary}

\begin{proof}
Part~(i): $c_{\max} = \max(1,1,1,2,2,2,2,2,3) = 3$ by inspection of
Table~\ref{tab:v5}.  The former maximum $c_{\max}=11$ is obsolete.

Part~(ii): Assemble the $N$-term sum left-to-right.  Each summand uses $\alpha_0$
nodes; each of the $N-1$ binary addition steps uses exactly 2 nodes (ADD-T1,
no domain restriction).  Total: $\alpha_0 N + 2(N-1)$.

Part~(iii): ADD-T1 is domain-unrestricted.  No case split on sign is needed.
\end{proof}

Table~\ref{tab:T40-comparison} quantifies the improvement for typical expressions.

\begin{table}[h]
\centering
\caption{T40 cost formula: comparison of v4 (add=11n) vs.\ v5 (add=2n).}
\label{tab:T40-comparison}
\smallskip
\begin{tabular}{@{}lcc@{}}
\toprule
\textbf{Expression ($N$ terms)} & \textbf{v4 cost} & \textbf{v5 cost} \\
\midrule
Shannon entropy $H=-\sum_{i=1}^N p_i\ln p_i$
  & $\alpha_0=3$: $14N-11$ & $5N-2$ \\
Taylor partial sum $\sum_{k=0}^{N-1} a_k x^k$
  & $\alpha_0=5$: $16N-11$ & $7N-2$ \\
Softmax numerator $\sum e^{x_i}$
  & $\alpha_0=1$: $12N-11$ & $3N-2$ \\
PK model $\sum A_i e^{-\lambda_i t}$
  & $\alpha_0=3$: $14N-11$ & $5N-2$ \\
\bottomrule
\multicolumn{3}{l}{\footnotesize
  $\alpha_0=\SB(\mulop)+\SB(\ln)=2+1=3$ for entropy terms;
  $\alpha_0=\SB(\powop)+\SB(\mulop)=3+2=5$ for polynomial terms (positive $x$).}
\end{tabular}
\end{table}

\begin{remark}[Interpretation]
The improvement from v4 to v5 is $9(N-1)$ nodes for every expression where
the dominant operation is summation with $\alpha_0=3$.  For large $N$, the saving
is roughly a factor of $14/5 \approx 2.8$ in total node count.  This directly
improves the compression ratio of SuperBEST v5 over the na\"ive baseline.
\end{remark}

%% =============================================================================
\section{Interaction with T01: Zero Finiteness Barrier}
\label{sec:T01}
%% =============================================================================

\subsection{Statement of T01}

\begin{theorem}[T01 --- Zero Finiteness]
Every finite $\F$-tree $T$ of $n$ nodes has at most $2^n - 1$ real zeros
on any bounded interval.  In particular, $T(x)=\sin(x)$ is impossible for
any finite $\F$-tree, since $\sin$ has infinitely many real zeros.
\end{theorem}

\subsection{Interaction: \Unchanged}

T01 is a theorem about the zero set of finite $\F$-trees, not about the cost
of any particular operation.  ADD-T1 does not change which functions are
$\F$-representable.

\begin{proposition}[ADD-T1 does not bypass T01]
\label{prop:T01-unchanged}
The following holds regardless of whether $\SB(\addop)=2$ or $\SB(\addop)=11$:
$\sin(x) \notin \ELC = \{\text{functions computable by finite }\F\text{-trees}\}$.
\end{proposition}

\begin{proof}
The zero of the tree $\LEdiv(x,\DEML(y,1))$ is the line $\{(x,y) : x=-y\}$,
which has measure 1 as a subset of $\R^2$.  More relevantly: the full 2-node
tree computing $x+y$ from inputs $(x,1)$ (treating $x$ as the sole variable)
has exactly one zero (at $x=-1$), consistent with T01 ($2^2-1 = 3 \geq 1$).
Building a Taylor partial sum $S_N(x) = \sum_{k=0}^{N-1} (-1)^k x^{2k+1}/(2k+1)!$
using cheap additions reduces the SuperBEST cost from $O(11N)$ to $O(9N-2)$,
but $S_N(x)$ is still a finite $\F$-tree, still subject to the zeros bound
$2^{\SB(S_N)}-1$.  Since $\sin$ has infinitely many zeros, no finite sum equals
$\sin$ exactly.
\end{proof}

\begin{corollary}[Cost of Taylor approximation to $\sin$]
\label{cor:sin-approx}
The $N$-term Taylor partial sum $S_N(x) = \sum_{k=0}^{N-1} (-1)^k x^{2k+1}/(2k+1)!$
satisfies:
\[
  \SB(S_N) \;\leq\; 9N - 2 \quad \text{(SuperBEST v5)}.
\]
The lower bound satisfies $\SB(S_N) \geq N$ (each of the $N$ distinct monomials
requires at least one node to distinguish it).  Hence $\SB(S_N) = \Theta(N)$.
\end{corollary}

\begin{proof}
Each term $(-1)^k x^{2k+1}/(2k+1)!$ has cost at most
$\SB(\powop)+\SB(\divop)+\SB(\negop) = 3+2+2 = 7$
(positive $x$; general $x$ may require additional handling for odd powers).
By T40 (Corollary~\ref{cor:T40-v5}), $N$ such terms sum to $7N+2(N-1)=9N-2$.
The lower bound follows because the $N$ monomials are algebraically independent
as subtrees.
\end{proof}

%% =============================================================================
\section{Interaction with the F16 Characterization Theorem ($\ELfin = \ELC$)}
\label{sec:ELC}
%% =============================================================================

\subsection{Statement}

The F16 Characterization Theorem states:
\[
  f:\R\to\R\text{ is exactly representable by a finite }\F\text{-tree}
  \;\iff\; f\in\ELC,
\]
where $\ELC$ is the exp-log closure: all functions obtainable by finitely many
applications of $\exp$, $\ln$, and rational arithmetic over $\Q$.

\subsection{Interaction: \Unchanged}

\begin{proposition}[ADD-T1 does not change $\ELC$]
$\ELC$ is a mathematical class of functions, not a cost class.
$\SB(\addop)=2$ is a statement about the node-count of the most efficient
$\F$-tree computing addition; it says nothing about which functions belong to $\ELC$.
\end{proposition}

\begin{proof}
$x+y$ was already known to lie in $\ELC$ (it is a rational arithmetic operation).
ADD-T1 shows it can be computed in 2 nodes rather than 11.  This reduces the cost
of a known member of $\ELC$; it does not introduce any new member.

Conversely, suppose some function $g\notin\ELC$ were suspected to benefit from
cheap addition.  If $g\notin\ELC$, then no finite $\F$-tree computes $g$ at any cost.
The cost of addition is irrelevant: even if $\SB(\addop)=0$, a function outside
$\ELC$ would remain outside $\ELC$.
\end{proof}

\begin{remark}[Cost classes vs.\ membership classes]
One may define the \emph{cost-$k$ class}
$\ELC_k = \{f : \SB(f) \leq k\}$.
ADD-T1 moves several functions from higher-cost classes to lower-cost classes:
for example, $x+y+z+w$ drops from $\ELC_{33}$ (v4, add=11n applied 3 times)
to $\ELC_6$ (v5, add=2n applied 3 times).
But $\ELC = \bigcup_{k\geq 0}\ELC_k$ is unchanged.
\end{remark}

%% =============================================================================
\section{Interaction with T36: Tree Counting Theorem}
\label{sec:T36}
%% =============================================================================

\subsection{Statement of T36}

T36 bounds the number of distinct $\F$-trees of a given node count.
Specifically, the number of binary $\F$-trees with exactly $N$ internal nodes
is at most $16^N \cdot C_N$, where $C_N = \binom{2N}{N}/(N+1)$ is the $N$-th
Catalan number.  T36 is used in pigeonhole arguments: if a target function $f$
imposes more than $16^k\cdot C_k$ independent algebraic constraints, then
$\SB(f)>k$.

\subsection{Interaction: \Unchanged (Qualitatively)}

T36 counts syntactically distinct trees; this count is independent of the cost
assigned to any operator.  However, the \emph{effective} role of T36 in the
proof system shifts modestly.

\begin{proposition}
The counting bound $16^N\cdot C_N$ is unchanged by ADD-T1.
However, the number of constraints required to distinguish $x+y$ from
nearby functions is unchanged (it remains $\geq 2$ independent constraints,
consistent with $\SB(\addop)\geq 2$).
\end{proposition}

\begin{remark}[ADD-T1 as a canonical 2-node subtree]
Under v5, $\addop$ has a canonical 2-node representation
$\LEdiv(x,\DEML(y,1))$.  This means that any $\F$-tree containing an
addition subexpression can now be rewritten with this specific 2-node subtree.
The fraction of all $N$-node trees that are ``equivalent'' (in the sense of
computing the same function) to a smaller tree may increase,
since addition targets are now achievable via a small canonical subtree.
This refinement does not change the statement of T36 but does affect the
efficiency of the SuperBEST optimizer.
\end{remark}

%% =============================================================================
\section{Interaction with T30: Depth Hierarchy}
\label{sec:T30}
%% =============================================================================

\subsection{Statement of T30}

T30 (EML Depth Spectrum Classification) establishes:
\begin{enumerate}
\item A strict hierarchy: $\EML_0 \subsetneq \EML_1 \subsetneq \EML_2 \subsetneq \cdots$
\item All standard elementary functions have EML depth $\leq 3$.
\item The canonical depth-$k$ witness is $\exp^{(k)}(x)$ (the $k$-fold iterated exponential).
\end{enumerate}

\subsection{Interaction: \Unchanged}

EML depth and SuperBEST cost are distinct but related metrics.  T30 classifies
functions by depth (the minimum over all $\F$-trees of the number of internal
nodes in the tree), which equals $\SB$ when no sharing or pattern bonuses apply.

\begin{proposition}[add achieves depth 2, consistent with T30]
The function $\addop(x,y)=x+y$ has EML depth equal to 2.
This is $\leq 3$, consistent with T30's census claim that all standard
elementary functions have depth $\leq 3$.
\end{proposition}

\begin{proof}
ADD-T1 exhibits a 2-node $\F$-tree for $\addop$; the 1-node lower bound is proved.
Hence EML depth of $\addop$ equals $\SB(\addop)=2$.
\end{proof}

\begin{remark}[Resolution of the ``depth census'']
Before ADD-T1, the SESSION-5 conjecture placed add at depth 11 (outside the
depth-$\leq 3$ census for standard elementary functions).  T30's claim that
all standard elementary functions achieve depth $\leq 3$ was therefore in tension
with the conjecture.  ADD-T1 resolves this tension completely: add sits at depth 2,
well within the T30 census.  The depth hierarchy ($\exp^{(k)}$ as canonical
depth-$k$ witness) is unaffected.
\end{remark}

%% =============================================================================
\section{Interaction with T35: Non-Mergeability Theorem}
\label{sec:T35}
%% =============================================================================

\subsection{Statement of T35}

T35 (Non-Mergeability) states that operator branches in a $\F$-tree cannot be
merged unless their output types are algebraically compatible.  More precisely,
an operator $O_1$ whose output type is $e^{(\text{linear form})}$ cannot be
directly combined via a second operator $O_2$ unless $O_2$ is defined to accept
inputs of that form.  T35 is used to rule out short trees by structural
incompatibility.

\subsection{Interaction: \Unchanged}

\begin{proposition}[ADD-T1 is T35-compatible]
The construction $\LEdiv(x,\DEML(y,1))$ satisfies the T35 compatibility conditions:
\begin{enumerate}
\item $\DEML(y,1)$ outputs $e^{-y}-\ln(1) = e^{-y}$, which is always in $(0,\infty)$.
\item $\LEdiv(x,z) = x - \ln z$ requires $z>0$.  Since $z = e^{-y} > 0$ for all
  $y\in\R$, the compatibility condition is met.
\end{enumerate}
Hence T35 presents no obstacle to the ADD-T1 construction.
\end{proposition}

\begin{remark}[T35 reinforces the 1-node lower bound]
T35 also provides an independent structural argument for why 1-node addition
is impossible: the output type of $\addop(x,y)=x+y$ is a bilinear function
with outputs ranging over all of $\R$, while no single $\F$-operator achieves
this output structure (all $\F$-operators involve $\exp$ or $\ln$, imposing
analytic constraints on their output range or derivative structure).
This is precisely what the derivative-obstruction proof in \texttt{AddLowerBound.lean}
formalises.
\end{remark}

%% =============================================================================
\section{Interaction with T37: exp/ln Disjointness}
\label{sec:T37}
%% =============================================================================

\subsection{Statement of T37}

T37 (Disjointness) states that the value sets (images) of distinct classes of
$\F$-operators are disjoint on relevant domains.  Specifically:
operators whose output is of the form $e^{(\cdot)}$ have positive-only images,
while operators of the form $\ln(\cdot)-(\cdot)$ can produce arbitrary real outputs.
A function whose value set requires both positive and negative outputs from
structurally incompatible branches cannot be computed in a single node.

\subsection{Interaction: \Unchanged}

\begin{proposition}[T37 confirms 1-node add is impossible]
The function $\addop(x,y)=x+y$ takes values in all of $\R$.  Every single
$\F$-operator either (a) is restricted to positive outputs (Class~I operators
of the form $e^{(\cdot)}$), or (b) has a specific analytic form (e.g., $\LEdiv$
is $x-\ln y$, which equals $x+y$ only if $\ln y = -y$ identically, which fails).
Hence T37 confirms $\SB(\addop)\geq 2$.
\end{proposition}

\begin{remark}[T37 and the 2-node solution]
The 2-node ADD-T1 construction is consistent with T37: $\DEML(y,1)=e^{-y}$
(positive, as required for the $\ln$ argument in $\LEdiv$), and the final
output $\LEdiv(x,e^{-y}) = x+y$ spans all of $\R$.
The two operators together ``span the necessary algebraic territory'' that no
single operator can reach.
\end{remark}

%% =============================================================================
\section{Interaction with the $\tan(1)$ Obstruction}
\label{sec:tan1}
%% =============================================================================

\subsection{Statement}

The $\tan(1)$ obstruction states:
$\tan(1) \notin \ELC$, i.e., $\tan(1)$ is not computable by any finite $\F$-tree.
This follows from the Lindemann--Weierstrass theorem: $\tan(1)$ is transcendental
over $\Q$, and $\ELC$ functions with rational parameters cannot produce this value.
Formally (DOOR-4): the contrapositive of the conditional chain shows
$\tan(1)\notin\EML_1$ is an unconditional theorem.

\subsection{Interaction: \Unchanged}

\begin{proposition}[Cheap add cannot produce $\tan(1)$]
$\SB(\addop)=2$ has no effect on the $\tan(1)$ obstruction.
$\tan(1)$ remains outside $\ELC$.
\end{proposition}

\begin{proof}
The transcendence argument is independent of how efficiently any specific
operation can be computed within $\ELC$.  Even if we could compute addition
in 0 nodes, the class $\ELC$ would remain the same (it is defined by finitely
many applications of $\exp$, $\ln$, rational arithmetic, which already contains
addition).  What matters is whether $\tan(1)$ can be expressed as a value of
some finite such composition --- and it cannot, by Lindemann--Weierstrass.
\end{proof}

\begin{remark}[Cost of approximations to $\tan$]
While $\tan(x)$ itself is not in $\ELC$, its Taylor partial sums are.
The $N$-term Maclaurin series
$T_N(x) = \sum_{k=0}^{N-1} B_{2k+2}\,4^{k+1}(4^{k+1}-1)\,x^{2k+1}/(2k+1)!$
has $\SB(T_N) \leq 9N-2$ under v5 (by Corollary~\ref{cor:sin-approx}, since
the structure is the same as for $\sin$).  But $T_N(1) \neq \tan(1)$ for
all finite $N$: the Taylor series converges but no finite partial sum equals
the limit.
\end{remark}

%% =============================================================================
\section{New Lemmas Arising from ADD-T1}
\label{sec:new-lemmas}
%% =============================================================================

The ADD-T1 breakthrough suggests three new formal lemmas that generalize
the interaction analysis.

\begin{lemma}[Signed-sum cost uniformity]
\label{lem:signed-sum}
Under SuperBEST v5, the cost of an $N$-term sum is independent of the signs
of the summands.  For any $f_1,\ldots,f_N \in \ELC$ each of SuperBEST cost $\alpha_0$:
\[
  \Cost\!\bigl(f_1 + f_2 + \cdots + f_N\bigr)
  \;=\; (\alpha_0 + 2)\,N - 2,
\]
whether the $f_i$ are positive, negative, or mixed.
\end{lemma}

\begin{proof}
ADD-T1 provides $\LEdiv(x,\DEML(y,1))=x+y$ for all $x,y\in\R$ without domain restriction.
Each binary addition step costs exactly 2 nodes regardless of the sign of the inputs.
Applying T40 (Corollary~\ref{cor:T40-v5}) gives the formula.
\end{proof}

\begin{lemma}[Universal binary upper bound]
\label{lem:universal-upper}
Every binary arithmetic operation $f$ in the SuperBEST v5 table satisfies:
\[
  \SB(f) \;\leq\; 3.
\]
The bound is tight: $\SB(\powop) = 3$.
\end{lemma}

\begin{proof}
Direct inspection of Table~\ref{tab:v5}: $\max\{\SB(f) : f\in\text{Table~\ref{tab:v5}}\} = 3$
(attained by $\powop$).  Before ADD-T1, the bound was 11.
\end{proof}

\begin{lemma}[Preliminary universal binary lower bound]
\label{lem:universal-lower}
For any binary function $f:\R^2\to\R$ that is not a function of a single variable
(i.e., $\partial f/\partial x \not\equiv 0$ and $\partial f/\partial y \not\equiv 0$),
if $f\in\ELC$ then $\SB(f) \geq 2$.
\end{lemma}

\begin{proof}[Proof sketch]
A 1-node $\F$-tree for $f$ would be a single $\F$-operator $O$ with
$O(x,y)=f(x,y)$.  The 16 operators of $\F$ are enumerable.  For each
operator $O$, either:
(a) $O$ does not depend on one of its inputs (e.g., $O(x,y)=O(x,c)$ for
constant $c$), so it cannot satisfy $\partial f/\partial y \neq 0$; or
(b) $O$ has the wrong functional form to match the partial derivative structure of $f$.
For the specific case of addition, the Lean proof in \texttt{AddLowerBound.lean}
provides concrete witnesses for all 16 operators.
A general proof for all two-variable ELC functions would require additional
algebraic-independence arguments; this is left as an open problem.
\end{proof}

\begin{remark}[Towards a general SB lower bound theorem]
Lemma~\ref{lem:universal-lower} is currently a theorem for specific operations
(add, sub, mul, neg, div in the SuperBEST table) and a conjecture for general
ELC operations.  A uniform proof would require showing that for any bivariate
$f\in\ELC$ with $\partial f/\partial x \not\equiv 0$ and $\partial f/\partial y\not\equiv 0$,
no single $\F$-operator matches $f$'s derivative structure.  The derivative-obstruction
technique of T08 / \texttt{AddLowerBound.lean} is the right tool but needs to be
cast in full generality.
\end{remark}

%% =============================================================================
\section{A Potential Universal Binary Cost Theorem}
\label{sec:universal}
%% =============================================================================

The analysis above raises a natural question:

\begin{definition}[Binary ELC operations]
A \emph{binary ELC operation} is any function $f:\R^2\to\R$ (or restricted to a
natural domain) that lies in $\ELC$.
\end{definition}

\begin{proposition}[Binary cost band]
\label{prop:band}
For every binary ELC operation $f$ in the SuperBEST v5 table:
\[
  2 \;\leq\; \SB(f) \;\leq\; 3.
\]
The lower bound $\SB(f)\geq 2$ holds for all operations that genuinely depend
on both inputs.  The upper bound $\SB(f)\leq 3$ holds by
Lemma~\ref{lem:universal-upper}.
\end{proposition}

\begin{proof}
Lower bound: every operation in the v5 table (Table~\ref{tab:v5}) other than
$\addop$, $\subop$, $\mulop$, $\divop$, $\negop$, $\recip$ has a structural
lower bound of $\geq 2$ by T08 (for the cases involving two independent variables)
or by the derivative obstruction.
Upper bound: Lemma~\ref{lem:universal-upper}.
\end{proof}

\begin{remark}[The conjecture for all binary ELC operations]
Proposition~\ref{prop:band} covers the 10 operations in the v5 table.
A stronger conjecture would be: \emph{for every two-variable ELC operation
$f$ with $\partial f/\partial x \not\equiv 0$ and $\partial f/\partial y\not\equiv 0$,
we have $2\leq\SB(f)\leq 6$.}
The upper bound of 6 is a very conservative estimate; in practice, all known
binary ELC operations achieve $\SB(f)\leq 3$.
This conjecture is consistent with all known data but remains open in full generality.
\end{remark}

%% =============================================================================
\section{Summary}
\label{sec:summary}
%% =============================================================================

Table~\ref{tab:summary} summarizes the interaction of $\SB(\addop)=2$ with
each proved structural result.

\begin{longtable}{@{}p{2.5cm}p{5.5cm}p{4.5cm}@{}}
\caption{Interaction of ADD-T1 ($\SB(\addop)=2$) with proved structural results.}
\label{tab:summary}\\
\toprule
\textbf{Theorem} & \textbf{Interaction} & \textbf{Net Effect} \\
\midrule
\endfirsthead
\multicolumn{3}{l}{\footnotesize\textit{(continued)}}\\
\toprule
\textbf{Theorem} & \textbf{Interaction} & \textbf{Net Effect} \\
\midrule
\endhead
\bottomrule
\endlastfoot
T08
(add$_{\text{pos}} \geq 2n$)
&
T08 proved $\SB(\addop_{\text{pos}})\geq 2$.
ADD-T1 proves $\SB(\addop)\leq 2$ for all reals.
Combined: $\SB(\addop)=2$ is a theorem.
Positive/general split eliminated.
&
\textbf{\textcolor{strengthened}{STRENGTHENED}}\\[6pt]
T40
(Linear Cost Law)
&
$c_{\max}$ drops from 11 to 3.
Formula becomes $(\alpha_0+2)N-2$.
Signed/unsigned sums unified.
&
\textbf{\textcolor{strengthened}{STRENGTHENED}}\\[6pt]
T01
(Zero Finiteness)
&
$\sin$, $\cos$, $\erf$ still outside $\ELC$.
Cheap add reduces Taylor approximation cost
to $9N-2$ (v5) but does not enable exact
representation.
&
\textbf{\textcolor{unchanged}{UNCHANGED}}\\[6pt]
$\ELfin = \ELC$
(Characterization)
&
$\ELC$ is a function class, not a cost class.
No new functions enter or leave $\ELC$.
Cost-$k$ classes $\ELC_k$ shift.
&
\textbf{\textcolor{unchanged}{UNCHANGED}}\\[6pt]
T36
(Tree Counting)
&
Tree counts are syntactic; unaffected by costs.
Add now has a canonical 2-node subtree.
&
\textbf{\textcolor{unchanged}{UNCHANGED}}\\[6pt]
T30
(Depth Hierarchy)
&
$\addop$ achieves depth 2, within T30's
depth-$\leq 3$ census.
Depth hierarchy unaffected.
&
\textbf{\textcolor{unchanged}{UNCHANGED}}\\[6pt]
T35
(Non-Mergeability)
&
ADD-T1 construction is T35-compatible.
T35 reinforces 1-node lower bound.
&
\textbf{\textcolor{unchanged}{UNCHANGED}}\\[6pt]
T37
(exp/ln Disjointness)
&
T37 confirms 1-node add impossible.
ADD-T1 shows 2 nodes span the required range.
&
\textbf{\textcolor{unchanged}{UNCHANGED}}\\[6pt]
$\tan(1)$ obstruction
&
Transcendence is independent of cost.
$\tan(1)\notin\ELC$ regardless of $\SB(\addop)$.
Taylor approximation costs $9N-2$.
&
\textbf{\textcolor{unchanged}{UNCHANGED}}\\
\end{longtable}

\subsection{New Lemmas (Summary)}

\begin{enumerate}
\item \textbf{Lemma~\ref{lem:signed-sum} (Signed-Sum Uniformity):}
  All $N$-term sums cost $(\alpha_0+2)N-2$ regardless of sign.

\item \textbf{Lemma~\ref{lem:universal-upper} (Universal Binary Upper Bound):}
  $\SB(f)\leq 3$ for all binary ELC operations in the v5 table.

\item \textbf{Lemma~\ref{lem:universal-lower} (Preliminary Binary Lower Bound):}
  $\SB(f)\geq 2$ for all two-variable ELC operations with non-trivial partial derivatives;
  fully proved for the v5 table operations, open in full generality.
\end{enumerate}

\subsection{Open Questions}

\begin{enumerate}
\item Can Lemma~\ref{lem:universal-lower} be proved for all binary ELC operations
  (not just the ten in the v5 table)?

\item Is $\SB(f)\leq 3$ for \emph{all} binary ELC operations,
  or does some pathological binary ELC function require 4 or more nodes?

\item Can Lemma~\ref{lem:universal-lower} be Lean-verified for the full
  16-operator family in a general-position statement?

\item Does $\SB(\addop)=2$ imply any improvement to the known bounds on
  the monogate attractor constant or related transcendence conjectures?
\end{enumerate}

%% =============================================================================
\section*{References}
%% =============================================================================

\begin{itemize}
\item \texttt{add\_gen\_2n\_proof.json} ---
  ADD-T1 primary record; construction and lower bound.
\item \texttt{lean/MonogateEML/MonogateEML/AddLowerBound.lean} ---
  Lean~4 sorry-free proof of $\SB(\addop)\geq 2$.
\item \texttt{python/paper/theorems/ADD\_T1\_General\_Addition\_2n.tex} ---
  companion theorem paper for ADD-T1.
\item \texttt{python/paper/theorems/SuperBEST\_v5\_Structural\_Audit.tex} ---
  complete v5 routing table with all lower bounds.
\item \texttt{python/paper/cost\_theory/R5\_Linear\_Cost\_Law.tex} ---
  T40 (Linear Cost Law); updated $c_{\max}=3$.
\item \texttt{python/paper/cost\_theory/R16\_T08\_Structural\_Proofs.tex} ---
  structural proofs for T08 entries (recip, pow, add).
\item \texttt{python/paper/theorems/F16\_Characterization.tex} ---
  $\ELfin = \ELC$ characterization theorem.
\item \texttt{python/paper/theorems/Depth\_Spectrum\_Classification.tex} ---
  T30 depth hierarchy.
\item \texttt{python/paper/theorems/Unified\_Structural\_Foundation.tex} ---
  T35, T36, T37, T40 techniques.
\item \texttt{python/paper/theorems/DOOR4\_Conditional\_Tan1.tex} ---
  $\tan(1)\notin\ELC$ obstruction.
\item \texttt{python/paper/proof\_t01\_lean\_closure.tex} ---
  T01 zero finiteness barrier.
\item \texttt{python/results/f17\_resolution.json} ---
  dual-domain correction (X20); positive vs.\ general domain table.
\item \texttt{python/results/add\_gen\_11n\_conjecture.json} ---
  SESSION-5 conjecture (superseded by ADD-T1).
\end{itemize}

\end{document}
